% 2_problema.tex

% Realice un análisis de la situación descrita y los desafíos tecnológicos que deben resolver para poder cumplir con el objetivo. Utilicen lenguaje técnico y preciso.
\begin{comment}
\section{Análisis del problema}

El objetivo del proyecto consiste en diseñar e implementar un sistema de comunicaciones digitales capaz de capturar, almacenar, transmitir y reproducir información de audio mediante un enlace inalámbrico de radiofrecuencia. La solución debe estar conformada por un dispositivo transmisor y un dispositivo receptor, ambos energéticamente autónomos y capaces de intercambiar información sin utilizar protocolos inalámbricos de alto nivel como WiFi o Bluetooth \cite{proyecto_ie0527}.

Para comprender mejor la arquitectura propuesta para este sistema de comunicaciones, en la Figura \ref{fig:diagrama_bloques2} se presenta un diagrama de bloques con la distintas etapas que conforman al sistema de comunicaciones.

\begin{figure}[H]
    \centering
    \includegraphics[width=1\linewidth]{images/Diagrama_audio2.png}
    \caption{Diagrama de bloques para la propuesta del sistema de comunicaciones.}
    \label{fig:diagrama_bloques2}
\end{figure}

La primera etapa corresponde a la captura de audio mediante un micrófono, el cual actúa como transductor electroacústico convirtiendo la voz humana en una señal eléctrica analógica, la cual corresponde a la señal de entrada del sistema.

Posteriormente, la señal analógica es procesada por un convertidor analógico-digital. Para la digitalización de voz se debe seleccionar adecuadamente la frecuencia de muestreo. Considerando que la mayor parte de la energía espectral de la voz inteligible se encuentra por debajo de los 3400 Hz, el Teorema de Nyquist establece que una frecuencia de muestreo de 8000 muestras por segundo es suficiente para reconstruir adecuadamente la señal \cite{lathi2009}.

En cuanto a la cuantización, se propone utilizar una resolución de 8 bits por muestra. Esta configuración permite representar cada muestra mediante 256 niveles discretos y proporciona una relación señal a ruido de cuantización aproximada de 49 dB, valor suficiente para aplicaciones de transmisión de voz \cite{lathi2009}.

Bajo estas condiciones, cada segundo de audio genera 8000 muestras, y cada muestra requiere 8 bits para su almacenamiento. Por consiguiente, la tasa de generación de datos resulta:

\begin{equation}
    R_b=f_sN = 8000 \cdot 8 = 64 \frac{\text{kbit}}{\text{s}}
\end{equation}

A partir de este dato se puede estimar cuál va a ser la cantidad de Bytes necesaria para trasmitir un mensaje de 5, 10 y 15 s.

\begin{equation}
\begin{aligned}
B_{5s} &= \frac{64000 \cdot 5}{8}
       = 40\;\text{kB}
\\[8pt]
B_{10s} &= \frac{64000 \cdot 10}{8}
        =80\;\text{kB}
\\[8pt]
B_{15s} & = \frac{64000 \cdot 15}{8}
        = 120\;\text{kB}
\end{aligned}
\label{kbytes}
\end{equation}

Dado que en este proyecto se requiere almacenar temporalmente el audio antes de su transmisión, el sistema debe contar con la memoria suficiente para almacenar mensajes con la duración más prolongada durante las pruebas \cite{proyecto_ie0527}. Por lo que una vez terminado el proceso de digitalización de la señal de audio, esta información debe almacenarse en un buffer de memoria. Una vez almacenada en memoria, la información digital debe prepararse para su transmisión inalámbrica, para facilitar el proceso de transmisión, se implementa un proceso de fragmentación, mediante el cual la secuencia de datos es dividida en paquetes más pequeños que puedan ser transmitidos individualmente.

La selección de la tecnología inalámbrica constituye otro de los principales desafíos del proyecto. Las especificaciones establecen que únicamente pueden emplearse bandas no licenciadas del espectro electromagnético y que no es permitido utilizar tecnologías de comunicación de alto nivel como WiFi, Bluetooth, Zigbee o LoRaWAN. En consecuencia, la comunicación debe implementarse directamente mediante módulos de radiofrecuencia de bajo nivel \cite{proyecto_ie0527}.

Se propone utilizar módulos de radiofrecuencia en la banda ISM de 2.4 GHz, dado que este permite transmitir hasta 2 Mbps, suficiente para transmitir los 120 kB en del mensaje más largo en menos de un segundo de transmisión.

Después de la transmisión inalámbrica, en el receptor se debe hacer el proceso inverso obtener una señal analógica a partir de la compresión realizada.

Entonces, una vez recibidos los paquetes, el sistema debe realizar un proceso de reensamblado que permita reconstruir los datos originales. Esta etapa también representa un riesgo de diseño, puesto que que debe de considerar diferentes escenarios, en los cuales paquetes pueden llegar en momentos diferentes, perder información durante la transmisión por el medio. Es por esto que resulta necesario implementar mecanismos de sincronización que permitan identificar el instante en el que se inicia una nueva transmisión, verificar que los datos se hayan recibido de forma correcta, para generar la señal digital reconstruida.

Posteriormene, esta señal digital reconstruida en el receptor debe ser almacenada momentáneamente para luego pasar por un proceso de conversión digital-analógico, con el fin de reconstruir la señal de audio pactada inicialmente, por último se debe colocar un parlante de tal forma que permita la reproducción audible del mensaje recibido en el receptor.

A partir de este análisis, se identifican cuatro desafíos tecnológicos principales que se pueden presentar durante el proceso de implementación, los cuales se en listan a continuación:

\begin{itemize}
    \item Fragmentar el audio en paquetes pequeños y reensamblarlos correctamente en el receptor.
    \item Sincronizar el receptor para que detecte el inicio de transmisión de forma autónoma.
    \item Manejar la pérdida de paquetes para evitar artefactos en el audio.
    \item Mantener la latencia total por debajo de los 60 segundos incluso para los mensajes más largos.
\end{itemize}
\end{comment}

\section{Análisis del problema}

El objetivo del proyecto consiste en desarrollar un sistema capaz de capturar, transmitir y reproducir información de audio de forma inalámbrica.
En esta sección, se analizan los principales desafíos que surgen a partir de los requerimientos solicitados en el enunciado.

Uno de los principales desafíos corresponde a la cadena completa de procesamiento de la señal.
Inicialmente, el transmisor debe convertir una señal de voz capturada mediante un micrófono en información digital para su posterior almacenamiento y transmisión.
El receptor debe reconstruir dicha información y reproducir el mensaje de audio.
Por esta razón, es necesario diseñar adecuadamente las etapas de captura, digitalización, almacenamiento, transmisión, recepción y reproducción, con el fin de minimizar las pérdidas de información y mantener una calidad de audio adecuada.
Dado que los módulos de radiofrecuencia de bajo nivel imponen un tamaño máximo de paquete del orden de decenas de bytes, el audio digitalizado debe fragmentarse en múltiples paquetes numerados, lo que exige diseñar un protocolo propio de segmentación y reensamblado en el receptor.

Otra limitación importante corresponde al método de comunicación utilizado.
Debido a que no se permite el uso de tecnologías inalámbricas de alto nivel como WiFi, Bluetooth o Zigbee, no se dispone de protocolos de comunicación previamente implementados.
En consecuencia, es necesario desarrollar un mecanismo propio de comunicación digital mediante módulos de radiofrecuencia de bajo nivel, lo que implica definir procedimientos de sincronización, transmisión y recepción de datos.

Adicionalmente, el receptor debe permanecer en estado de espera de forma permanente y detectar autónomamente el inicio de una transmisión válida.
Esto representa un reto de sincronización, debido a que el receptor debe diferenciar un paquete de inicio de transmisión de posibles señales del ambiente o ruido presente en el canal, sin conocer previamente el instante de llegada de la información.

También se presenta una restricción asociada al uso de un único microprocesador por dispositivo, pues todas las tareas relacionadas con captura de audio, almacenamiento, procesamiento digital, control del enlace inalámbrico y reproducción deben ejecutarse mediante una sola unidad de procesamiento.
Esto obliga a optimizar el uso de los recursos computacionales disponibles.

Otro aspecto relevante corresponde al almacenamiento de la información digitalizada.
El sistema debe ser capaz de almacenar temporalmente mensajes de voz de hasta $15$ segundos de duración antes de su transmisión.
La cantidad de memoria requerida depende directamente de parámetros como la frecuencia de muestreo, el número de bits por muestra y el posible uso de técnicas de compresión.
Por tanto, es necesario estimar la capacidad de almacenamiento requerida para garantizar que todos los mensajes puedan ser procesados sin inconvenientes.

De forma complementaria, la cantidad de información generada durante la digitalización afecta directamente el tiempo necesario para completar la transmisión.
Por esta razón, resulta conveniente analizar la implementación de mecanismos de compresión de audio que permitan reducir el tamaño de los datos transmitidos.
La aplicación de técnicas de compresión puede disminuir el tiempo de transmisión y los requerimientos de almacenamiento, aunque introduce el desafío adicional de mantener una calidad de audio suficiente para garantizar la inteligibilidad del mensaje recibido.

Con respecto al desempeño del sistema, este debe cumplir simultáneamente requisitos de tiempo de respuesta y confiabilidad.
El proceso completo de captura, procesamiento, transmisión y reproducción debe finalizar dentro de los límites temporales establecidos en los requerimientos del proyecto.
Además, se debe garantizar que el mensaje reproducido represente adecuadamente la información capturada por el micrófono.

La transmisión inalámbrica está expuesta a fenómenos de ruido, interferencia y pérdida de paquetes en el canal, lo que puede afectar la calidad del audio reconstruido.
También, la variación temporal en la llegada de paquetes (\textit{jitter}) puede producir discontinuidades en la reproducción del audio si las muestras no se entregan a la tasa correcta, por lo que se requiere un mecanismo de almacenamiento temporal en el receptor que absorba estas variaciones.
Por estas razones, resulta necesario considerar mecanismos que permitan verificar la integridad de la información recibida y aumentar la confiabilidad del sistema.

Otro aspecto a considerar es la autonomía energética.
Ambos dispositivos deben operar mediante baterías y mantener un consumo reducido durante su funcionamiento.
Esto obliga a considerar el consumo asociado a los módulos de procesamiento, almacenamiento, captura de audio, transmisión por radiofrecuencia y reproducción, con el fin de garantizar una autonomía adecuada para la operación del sistema.

En resumen, a partir del análisis realizado, el principal desafío consiste en satisfacer simultáneamente los requisitos de calidad de audio, latencia máxima de $60$ segundos por mensaje y confiabilidad del enlace, bajo las restricciones de ancho de banda del canal, tamaño de paquete del módulo RF y autonomía energética de los dispositivos.